\chapter{Scope, vocabulary and method}\label{ch:scope}

\section{What was reviewed, and against what}

The object of the review is the \emph{blueprint} of the leanVM proof system in leanerVM, the
document \file{docs/roadmap/protocol-blueprint.md}, together with the code already built from it.
leanerVM was reviewed on its branch \code{main} at commit \code{b435631} (local and remote agree).
\Cref{tab:revisions} lists every revision a statement of this report can be about. A statement
about a library is about its \emph{pin}, the revision leanerVM builds against, unless it says
\enquote{upstream}.

\begin{table}[h]
\centering\small
\begin{tabular}{@{}llll@{}}
\toprule
What & Revision & Where it was read & Role \\
\midrule
leanerVM & \code{b435631} (\code{main}) & the working directory & the object of the review \\
leanVM specification & \code{a386121f} & \file{doc/leanvm/body/*.tex} & ground truth: the intended protocol \\
leanVM Rust verifier & \code{a386121f} & \file{crates/} & ground truth: the deployed protocol \\
leanVM Python verifier & \code{a386121f} & \file{python-verifier/verifier.py} & ground truth: second implementation \\
ArkLib & \code{dca90385} & \file{.lake/packages/Arklib} & the framework of the statements \\
VCVio & \code{f9dc47d9} & \file{.lake/packages/VCVio} & probabilistic oracle computations \\
CompPoly & \code{3468b38c} & \file{.lake/packages/CompPoly} & fields and computable polynomials \\
Clean & \code{93c9d1ef} & \file{.lake/packages/Clean} & the circuit language of the arithmetization \\
Mathlib, Lean & \code{v4.33.1} & \file{.lake/packages/mathlib} & mathematics, and the kernel \\
\bottomrule
\end{tabular}
\caption{The revisions this report speaks about.}\label{tab:revisions}
\end{table}

Three sources define leanVM at the pin, and they do not always agree. This report uses two words
for them and keeps them apart:
\begin{itemize}
\item the \textbf{specification} is the tex source at the pin: what leanVM is meant to be;
\item the \textbf{deployed} protocol is what the Rust verifier at the pin accepts: what would run
  on Ethereum. The Python verifier is a second implementation of it.
\end{itemize}
A disagreement between them is a finding in its own right (\cref{ch:faithfulness}), and every
theorem discussed is labelled with the verifier it describes. A PDF of the specification
circulates (\file{leanVM-b-2.pdf}); it was built from an earlier revision, differs from the
pinned text, and is never cited here.

\section{The three criteria}

\begin{description}
\item[Faithfulness.] The blueprint models leanVM's proof system interaction by interaction: every
  prover message, every verifier challenge and every verifier check, in leanVM's order.
  A divergence is classified as one of: \emph{an error of the blueprint}, to be fixed;
  \emph{a deliberate deviation}, to be kept with its reason recorded and with the theorem that
  relates it to the deployed verifier named and owed; or \emph{a deviation forced by an upstream
  library}, recorded with its workaround and the condition under which it is retired.
\item[Non-vacuity.] A definition is inhabited by the object leanVM intends and refuted by the
  near misses; a hypothesis can be met, and by the real protocol, not only by a trivial one; an
  error bound means something at leanVM's parameters; an extractor computes. A statement that
  typechecks and a proof that compiles are evidence of neither.
\item[Auditability.] The declarations a reader has to read and trust in order to know what the
  master theorems say are as few and as short as they can be.
\end{description}
Faithfulness and non-vacuity of the relation, the seams and the master theorems come first;
the size of the audit surface comes second and can be improved iteratively once the statements
are right.

\section{Vocabulary}\label{sec:vocabulary}

The repository's documents use several names for overlapping things. This report uses one name
for each, as follows; \cref{ch:documentation} proposes that the documents do the same.

\begin{longtable}{@{}p{0.2\linewidth}p{0.76\linewidth}@{}}
\toprule
Term & Meaning \\
\midrule
\endhead
blueprint & \file{docs/roadmap/protocol-blueprint.md}, the specification of what is wanted for
  the proof system. It titles itself \enquote{Roadmap}; the two are one document. \\
status & \file{docs/roadmap/protocol-status.md}. \\
tracker & issue \#12 of the repository, with the long comment that holds the specifications of
  the holes (the \emph{hole comment}). \\
oracle protocol & the interactive protocol in the ideal model where the committed polynomial is
  an oracle the verifier may query and challenges are uniformly random. In the literature: a
  polynomial interactive oracle proof. \\
phase & one step of the oracle protocol: commit, bus, table sumcheck, public input, Flock,
  opening. \\
spine & the part of the blueprint, and of the code, that fixes what the phases share: the
  abstract instance \lean{M3Instance}, the relation \lean{M3Holds}, the seam relations, the
  interfaces a phase has to meet, the composition and the two master theorems. \\
seam & the relation that holds between two consecutive phases: the output relation of one and,
  by definition, the input relation of the next. \\
hole & a hole the spine leaves open (a phase, a generic component, the adaptor, the
  compilation). \\
layer & the blueprint's numbering of the same work (Layers 0 to 13). \\
adaptor & the part, not yet built, that instantiates the abstract instance at leanISA and
  carries the proof system's relation to leanISA's constraint relation \lean{SatisfiedBy}. \\
compilation & what turns the oracle protocol into a non-interactive proof: the polynomial
  commitment (WHIR over Merkle trees) in place of the oracle, and Fiat--Shamir in place of the
  verifier's coins. \\
master theorems & \lean{piop_perfectCompleteness} and \lean{piop_rbrKnowledgeSoundness}, and the
  theorems the blueprint builds on them, not yet proved: \lean{verify_knowledgeSound} for the
  compiled verifier, and \lean{baseVerifier_extractsExecution} with \lean{baseProver_complete},
  the proof system's target theorem. \\
pin & the revision of a dependency recorded in \file{upstreams.json}. \\
load-bearing & said of a declaration that a master theorem's statement unfolds to, so that a
  reader must read it to know what the theorem says. \\
Category A, Category B & the blueprint's two kinds of content: A is what the verifier accepts
  and why it is sound, written from the specification and then diffed against the code; B is
  transcribed data (layouts, encodings, constants, parameters), copied from its source and
  wrong if it deviates. \\
Layer $N$, acceptance test $N$, decision $N$ & the blueprint's own numbering of its sections,
  its tests and its decisions; the report keeps these numbers, since they are the document's
  structure and not codes for other things. \\
\bottomrule
\end{longtable}

\paragraph{Letter codes.} The repository's documents label holes, ledger entries, findings and
decisions with a letter and a number. This report names things in words. Where a code helps to
find the thing in the repository it is given once, in parentheses, after the name.
\Cref{ch:codeindex} is the index of every code met. The report has one code family of its
own, the register's row numbers \code{R1} to \code{R145} (\cref{ch:register}), used only to
point at a finding beside its name.

\section{Method}\label{sec:method}

The review was run by one orchestrating session and independent sub-agents, each with its own
context and a self-contained task, all bound by the limits of the session: no edit to a tracked
file, nothing posted to an issue tracker, no checkout moved.

\begin{enumerate}
\item \textbf{Ground truth first.} For each part of the protocol a sub-agent read the
  specification, the Rust and the Python verifier at the pin and wrote down the transcript
  (who sends what, in which order, which challenge follows which message), every check of the
  verifier with its rejection path, and the errors claimed. Only then was the blueprint set
  against it.
\item \textbf{The code as evidence.} The built code (the spine, Layers 0 and 1, the public-input
  phase) was catalogued declaration by declaration from the source, classified as load-bearing,
  interface, helper or test fixture, and read adversarially against its docstrings and the
  blueprint.
\item \textbf{Probes.} Claims of non-vacuity were tested by Lean files that build against
  \code{main}: inhabitants, near misses, and mutations of a verifier that show where a proof
  stops when a check is removed. The probes are scratch work and are reproduced in
  \cref{ch:probes}.
\item \textbf{Validation of findings and of their absence.} Every finding carries its evidence:
  the line of the pinned source, the Lean declaration, and the probe or counterexample that
  reproduces it. Where a part was found to have no issue, the report says what was checked.
  Every citation of the four ground-truth dossiers was checked against the sources by a
  second agent, and eleven of the thirty-seven major findings (those not concerning the
  opening, whose dossier landed last) were re-derived by a hostile second reader before this
  report was written; the others rest on their dossier's evidence and the citation checks.
\end{enumerate}

\paragraph{Severity.} \sev{critical}: a theorem or definition is wrong, vacuous or unfaithful in
a way that would let an unsound verifier be \enquote{proved}. \sev{major}: the blueprint
specifies something leanVM does not do, or omits something it does, or a statement cannot be
proved or stated as written. \sev{minor}: imprecision, stale text, avoidable audit surface.
\sev{note}: worth recording.

\section{What this report is not}

It is the record of a review, not a further source of truth about the proof system. Whatever it
concludes should change is written as a proposed change to the blueprint (\cref{ch:changes}),
with the status and the tracker following, so that the blueprint remains the one place that
says what is wanted. Nothing proposed here has been applied.
